\documentclass[11pt,a4paper]{article}

\usepackage[utf8]{inputenc}

\usepackage[T1]{fontenc}

\usepackage[polish]{babel}

\usepackage{lmodern}

\usepackage{geometry}

\geometry{margin=27mm}

\usepackage[hidelinks]{hyperref}

\title{Rozdzielność bliskości i autonomii w różnych kontekstach}

\author{Mateusz Pomianek}

\date{Szkic naukowy, wersja 1.0.0 | 28 listopada 2026}

\begin{document}

\maketitle

\begin{abstract}

Proponowana miara główna to różnica predykcyjnego wyniku logarytmicznego modeli jedno- i dwuczynnikowego.

\end{abstract}

\section{Wprowadzenie}

Teoria autodeterminacji ujmuje autonomię i więź jako odrębne potrzeby (Patrick et al., 2007). Rozdzielność wymiarów jest tu testowana w ocenach relacji w konkretnych sytuacjach.

\section{Cel i pytania badawcze}

Hipotezy robocze wymagają zatwierdzenia przed rozpoczęciem właściwej weryfikacji.

\paragraph{H1.} W niezależnej próbie model z dwoma skorelowanymi czynnikami bliskości i autonomii uzyska lepszą trafność predykcyjną odpowiedzi niż model jednoczynnikowy.

\paragraph{H2.} Manipulacja dostępnością wsparcia podniesie ocenę bliskości bardziej niż ocenę autonomii, a manipulacja swobodą decyzji podniesie ocenę autonomii bardziej niż ocenę bliskości.

\section{Materiał i metody}

Pierwszorzędowy wskaźnik: różnica predykcyjnego wyniku logarytmicznego modeli jedno- i dwuczynnikowego.

W winietach należy niezależnie manipulować oboma czynnikami w kontekście konfliktu i codziennej decyzji. Trafność predykcyjna jest oceniana logarytmicznym wynikiem prognoz na osobach niewidzianych w dopasowaniu. Porównanie średnich latentnych między kontekstami wymaga co najmniej częściowej inwariancji progów i ładunków.

Przed badaniem z udziałem ludzi należy zatwierdzić wymagania etyczne, zasady zgody i ochronę danych.

\section{Plan analizy}

Zamrożenie estymandu, reguł oceny jakości i obsługi braków nastąpi przed dostępem do danych wynikowych. Raport obejmie oszacowanie kontrastu głównego z przedziałem niepewności i pełne mianowniki. Jeśli przewidziano hipotezy drugorzędowe, przed analizą należy ustalić sposób kontroli wielokrotności.

\section{Wyniki}

W niniejszym pliku nie umieszczono wyników empirycznych, statystyk ani nieprzeprowadzonych analiz. Sekcja wymaga uzupełnienia po uzyskaniu i weryfikacji materiału dowodowego.

\section{Dyskusja i ograniczenia}

Ograniczenia obejmują możliwe błędy pomiaru, dobór próby, alternatywne mechanizmy i ograniczoną możliwość uogólniania. Korelacja nie dowodzi przyczynowości; poprawność techniczna nie jest dowodem znaczenia klinicznego.

\section{Wnioski}

Hipotezy i sposób pomiaru stanowią podstawę oceny wykonalności oraz przygotowania właściwego protokołu.

\section*{Bibliografia}

{\setlength{\parindent}{0pt}
\setlength{\parskip}{5pt}
\hangindent=1.27cm
Patrick, H., Knee, C. R., Canevello, A., \& Lonsbary, C. (2007). The role of need fulfillment in relationship functioning and well-being: A self-determination theory perspective. \textit{Journal of Personality and Social Psychology, 92}(3), 434--457. \url{https://doi.org/10.1037/0022-3514.92.3.434}\par
}

\end{document}
